\documentclass[../../main.tex]{subfiles}
\begin{document}
% SOURCE p=236; printed=221; object=ch07-czero-dual-title; region=middle
\section{$C_0(X)$ 的对偶}\label{section:ch07-czero-dual}
% SOURCE p=236; printed=221; object=ch07-positive-czero-dual-introduction; region=middle-lower
回顾：对任意 LCH 空间 $X$，$C_0(X)$ 是 $C_c(X)$ 的一致闭包（\cref{proposition:ch04-cc-dense-czero}）。因此，若 $\mu$ 是 $X$ 上的 Radon 测度，泛函 $I(f)=\int f\,\mathrm d\mu$ 在关于一致范数有界时可以连续延拓到 $C_0(X)$。由
\begin{align*}
\mu(X)=\sup\left\{\int f\,\mathrm d\mu:f\in C_c(X),\ 0\leqslant f\leqslant1\right\}
\end{align*}
（\cref{align:ch07-riesz-open-sup}的特例），以及 $|\int f\,\mathrm d\mu|\leqslant\int|f|\,\mathrm d\mu$，可知这种情形恰在 $\mu(X)<\infty$ 时发生，此时 $\mu(X)$ 就是 $I$ 的算子范数。

所以，$C_0(X)$ 上的正有界线性泛函已经得到刻画：它们就是关于有限 Radon 测度的积分。本节要把这个结果扩展成对 $C_0(X)^*$ 的完整描述。关键事实是：$C_0(X,\mathbb R)$ 上的实线性泛函具有一种“Jordan 分解”。
% SUPPLEMENT: id=v14-ch07-positive-czero-details; parent=ch07-positive-czero-dual-introduction; kind=proof-detail+motivation+hypotheses; origin=independent-derivation
\begin{remark}
\textbf{延拓、正性与范数的完整核验（重构新增）。}若 $\mu(X)<\infty$，则每个 $f\in C_0(X)$ 有界可积，且 $|I(f)|\leqslant\mu(X)\|f\|_u$。上面的开集公式又给出 $\|I\|\geqslant\mu(X)$，故范数相等。反过来，若 $I$ 在 $C_c(X)$ 的一致范数下有界，所有 $0\leqslant f\leqslant1$ 的测试积分均不超过 $\|I\|$，所以 $\mu(X)\leqslant\|I\|<\infty$。有界泛函由稠密性唯一延拓到 $C_0(X)$：对一致趋于 $f$ 的 $f_n\in C_c(X)$，取 $\lim I(f_n)$，范数界保证极限存在且与近似列无关；它与直接积分一致。

若先给定正有界泛函 $I\in C_0(X)^*$，限制到 $C_c(X)$ 后由 Riesz 定理产生 Radon 测度，范数界使其有限，上述连续延拓唯一性恢复原 $I$。反向，有限非负测度的积分显然保持非负性。动机是把“每个紧支集上的局部界”升级成“全空间的一致界”；新增的有限\emph{总}质量恰是所需条件，紧集有限本身不足以推出有界性。此处尚未预设所有正泛函自动有界，该事实将在\cref{exercise:ch07-positive-czero-automatic-bound}证明。
\end{remark}
% END SUPPLEMENT: id=v14-ch07-positive-czero-details

% SOURCE p=236; printed=221; object=ch07-functional-jordan; number=7.15; region=lower
\begin{lemma}\label{lemma:ch07-functional-jordan}
若 $I\in C_0(X,\mathbb R)^*$，则存在正泛函 $I^+,I^-\in C_0(X,\mathbb R)^*$，使 $I=I^+-I^-$。
\end{lemma}
% SOURCE p=236,237; printed=221-222; object=ch07-functional-jordan-proof; region=bottom-upper
\begin{proof}
对 $f\in C_0(X,[0,\infty))$，定义
\begin{align*}
I^+(f)=\sup\{I(g):g\in C_0(X,\mathbb R),\ 0\leqslant g\leqslant f\}.
\end{align*}
因为 $0\leqslant g\leqslant f$ 时 $|I(g)|\leqslant\|I\|\|g\|_u\leqslant\|I\|\|f\|_u$，且 $I(0)=0$，有 $0\leqslant I^+(f)\leqslant\|I\|\|f\|_u$。断言 $I^+$ 是某个线性泛函在非负锥上的限制；其证明与\cref{section:ch02-complex-integration}中积分线性的证明很相似。

显然当 $c\geqslant0$ 时 $I^+(cf)=cI^+(f)$。若 $0\leqslant g_1\leqslant f_1$、$0\leqslant g_2\leqslant f_2$，则 $0\leqslant g_1+g_2\leqslant f_1+f_2$，故 $I^+(f_1+f_2)\geqslant I(g_1)+I(g_2)$。取上确界得到 $I^+(f_1+f_2)\geqslant I^+(f_1)+I^+(f_2)$。另一方面，若 $0\leqslant g\leqslant f_1+f_2$，令 $g_1=\min(g,f_1)$、$g_2=g-g_1$，则 $0\leqslant g_1\leqslant f_1$、$0\leqslant g_2\leqslant f_2$。于是 $I(g)=I(g_1)+I(g_2)\leqslant I^+(f_1)+I^+(f_2)$；再取上确界，得反向不等式。所以 $I^+(f_1+f_2)=I^+(f_1)+I^+(f_2)$。

现对任意实值 $f\in C_0(X)$，其正负部分 $f^+,f^-$ 都在非负锥中，定义 $I^+(f)=I^+(f^+)-I^+(f^-)$。若也有 $f=g-h$，其中 $g,h\geqslant0$，则 $g+f^-=h+f^+$，所以 $I^+(g)+I^+(f^-)=I^+(h)+I^+(f^+)$。因此 $I^+(f)=I^+(g)-I^+(h)$，由此容易像\cref{proposition:ch02-real-integral-linearity}的证明那样得到 $I^+$ 在 $C_0(X,\mathbb R)$ 上线性。此外，
\begin{align*}
|I^+(f)|&\leqslant\max\{I^+(f^+),I^+(f^-)\}\\
&\leqslant\|I\|\max\{\|f^+\|_u,\|f^-\|_u\}=\|I\|\|f\|_u,
\end{align*}
故 $\|I^+\|\leqslant\|I\|$。最后，令 $I^-=I^+-I$，则它属于对偶空间，而由 $I^+$ 的定义，$I^+$ 与 $I^-$ 都为正泛函。
\end{proof}
% SUPPLEMENT: id=v14-ch07-functional-jordan-details; parent=ch07-functional-jordan; kind=proof-detail+motivation+hypotheses; origin=independent-derivation
\begin{remark}
\textbf{正锥到全空间的线性步骤（重构新增）。}候选零函数保证上确界非空，$I$ 的有界性保证它有限。$c>0$ 时 $g\mapsto cg$ 给出 $[0,f]$ 与 $[0,cf]$ 之间的双射，故正齐次；$c=0$ 时候选仅零函数，结论也成立。上确界加性第一向可分别取两个距上确界小于 $\varepsilon/2$ 的候选，求和后令 $\varepsilon\downarrow0$。第二向的 $\min(g,f_1)$ 连续，且其模不超过 $f_1$，故在无穷远处消失；$g_2=g-g_1$ 也属于 $C_0$，逐点分 $g\leqslant f_1$ 与 $g>f_1$ 可见 $g_2\leqslant f_2$。

把实函数写成任意非负差 $f=g-h$ 后，表示独立性已由原证建立。若 $f=g-h$、$u=v-w$，则 $f+u=(g+v)-(h+w)$，非负锥加性给出 $I^+(f+u)=I^+(f)+I^+(u)$；非负标量齐次由同样的差表示得到，负标量再用 $-f=h-g$ 得到。因此全空间的实线性完整成立。对 $f\geqslant0$，$g=f$ 是上确界中的候选，故 $I^+(f)\geqslant I(f)$，这正好给出 $I^-(f)\geqslant0$。有界性由 $\|I^-\|\leqslant\|I^+\|+\|I\|$ 得到。

动机与测度的 Jordan 分解相同：在正输入 $f$ 所允许的区间 $0\leqslant g\leqslant f$ 内，取 $I$ 能检测到的最大正贡献。实标量域提供顺序及正负分解；复泛函须先拆成实、虚两部分。$C_0$ 对正负部分、有限最大最小和线性组合封闭，是这项构造的关键结构。这里的上确界构造给出指定的 $I^+$，不表示任意两个正泛函之差的表示唯一；其与测度的真正 Jordan 部分的关系见\cref{exercise:ch07-functional-measure-jordan}。
\end{remark}
% END SUPPLEMENT: id=v14-ch07-functional-jordan-details

% SOURCE p=237; printed=222; object=ch07-complex-functional-four-measures; region=middle
任意 $I\in C_0(X)^*$ 都由它在 $C_0(X,\mathbb R)$ 上的限制 $J$ 唯一确定，而 $J=J_1+\mathrm iJ_2$，其中 $J_1,J_2$ 为实线性泛函。因此，由\cref{lemma:ch07-functional-jordan}及其前面的讨论，存在有限 Radon 测度 $\mu_1,\ldots,\mu_4$，使 $I(f)=\int f\,\mathrm d\mu$，其中 $\mu=\mu_1-\mu_2+\mathrm i(\mu_3-\mu_4)$。
% SUPPLEMENT: id=v14-ch07-complex-four-measures-details; parent=ch07-complex-functional-four-measures; kind=proof-detail+motivation+hypotheses; origin=independent-derivation
\begin{remark}
\textbf{从实限制恢复复泛函（重构新增）。}对实函数 $u$，令 $J_1(u)=\operatorname{Re}I(u)$、$J_2(u)=\operatorname{Im}I(u)$，它们实线性且 $|J_j(u)|\leqslant\|I\|\|u\|_u$。对每个 $J_j$ 用前引理拆为两个正有界泛函，再用有限正测度表示，得到四个测度。对实 $u$ 已有 $I(u)=\int u\,\mathrm d\mu$；任意复函数 $f=u+\mathrm iv$，其两部分都属于 $C_0(X,\mathbb R)$，复线性给出 $I(f)=I(u)+\mathrm iI(v)=\int f\,\mathrm d\mu$。四个有限正测度的组合是有限变差复测度。这个步骤证明表示的存在性；其 Radon 性及唯一性需下面的定义与定理，不能仅因写成四个正测度就把它们未经证明地认作 $\mu$ 的典范变差分量。
\end{remark}
% END SUPPLEMENT: id=v14-ch07-complex-four-measures-details

% SOURCE p=237; printed=222; object=ch07-signed-complex-radon-definition; region=middle-lower
\begin{definition}[符号与复 Radon 测度]\label{definition:ch07-signed-complex-radon}
\textbf{符号 Radon 测度}是正、负变差都为 Radon 测度的符号 Borel 测度。\textbf{复 Radon 测度}是实部和虚部都为符号 Radon 测度的复 Borel 测度。用 $M(X)$ 表示 $X$ 上的复 Radon 测度所成的空间；对 $\mu\in M(X)$，定义
\begin{align*}
\|\mu\|=|\mu|(X),
\end{align*}
其中 $|\mu|$ 为 $\mu$ 的全变差。
\end{definition}
在第二可数 LCH 空间上，每个复 Borel 测度都是 Radon 测度：复测度有界，所以这由\cref{theorem:ch07-automatic-regularity}推出。
% SUPPLEMENT: id=v14-ch07-complex-radon-definition-details; parent=ch07-signed-complex-radon-definition; kind=proof-detail+motivation+hypotheses; origin=independent-derivation
\begin{remark}
\textbf{动机与有限性约定（重构新增）。}对符号测度直接谈大小逼近会受到正负抵消影响，因此定义把正则性分别施加于 Jordan 的正、负变差；复测度再分实、虚两部分。这里复测度沿用第三章约定，在所有可测集上取有限复值，其全变差总质量有限；一般的符号 Radon 测度则不必总变差有限，不能自动把它都纳入 $M(X)$。$M(X)$ 的标量域是 $\mathbb C$，其“空间”和“范数”性质将在下一命题证明。

上面的第二可数断言具体如下：$\operatorname{Re}\mu$、$\operatorname{Im}\mu$ 的四个非负 Jordan 分量都有限，因而在紧集上有限；由自动正则定理各为 Radon 测度，恰好满足本定义。有限性本身在一般紧 Hausdorff 空间仍不足以保证正则，\cref{exercise:ch07-ordinal-measure}已给出反例。
\end{remark}
% END SUPPLEMENT: id=v14-ch07-complex-radon-definition-details

% SOURCE p=237; printed=222; object=ch07-variation-radon; number=7.16; region=lower
\begin{proposition}\label{proposition:ch07-variation-radon}
若 $\mu$ 为复 Borel 测度，则 $\mu$ 为 Radon 测度当且仅当 $|\mu|$ 为 Radon 测度。此外，$M(X)$ 是向量空间，$\mu\mapsto\|\mu\|$ 是其上的范数。
\end{proposition}
% SOURCE p=237; printed=222; object=ch07-variation-radon-proof; region=bottom
\begin{proof}
先注意：由\cref{proposition:ch07-sigma-finite-inner,proposition:ch07-open-closed-approximation}，有限正 Borel 测度 $\nu$ 为 Radon，当且仅当对每个 Borel 集 $E$ 和每个 $\varepsilon>0$，都存在紧集 $K$ 和开集 $U$，使 $K\subseteq E\subseteq U$ 且 $\nu(U\setminus K)<\varepsilon$。第一个断言容易由此推出。事实上，写 $\mu=\mu_1-\mu_2+\mathrm i(\mu_3-\mu_4)$。若 $|\mu|(U\setminus K)<\varepsilon$，则所有 $j$ 都有 $\mu_j(U\setminus K)<\varepsilon$；反过来，若 $\mu_j(U_j\setminus K_j)<\varepsilon/4$ 对所有 $j$ 成立，则令 $K=\bigcup_{j=1}^4K_j$、$U=\bigcap_{j=1}^4U_j$，可得 $|\mu|(U\setminus K)<\varepsilon$（原书此处印作 $\mu(U\setminus K)<\varepsilon$，全变差符号的订正见下文）。同样的论证表明 $M(X)$ 对加法和数乘封闭。最后，$\|\cdot\|$ 的范数性质由\cref{proposition:ch03-complex-variation-triangle}推出。
\end{proof}
% SUPPLEMENT: id=v14-ch07-variation-radon-details; parent=ch07-variation-radon; kind=proof-detail+correction+motivation+hypotheses; origin=independent-derivation
\begin{remark}
\textbf{夹层判据、典范分量与范数（重构新增）。}有限正 Radon 测度在每个 Borel 集上内、外正则，分别选误差小于 $\varepsilon/2$ 的紧内罩和开外罩即得夹层判据。反过来，该判据同时给出所有 Borel 集的内、外正则性，而紧集有限由总质量有限得到，故确为 Radon 测度。

第一向应取 $\mu_1=(\operatorname{Re}\mu)^+$、$\mu_2=(\operatorname{Re}\mu)^-$、$\mu_3=(\operatorname{Im}\mu)^+$、$\mu_4=(\operatorname{Im}\mu)^-$ 这四个\emph{典范}分量。它们满足 $\mu_j\leqslant|\mu|$：若 $\mathrm d\mu=h\,\mathrm d|\mu|$ 且 $|h|=1$ 几乎处处，则相应密度是 $(\operatorname{Re}h)^\pm$、$(\operatorname{Im}h)^\pm$，都在 $[0,1]$ 内。因此 $|\mu|$ 的小夹层同时是四个分量的小夹层。不能对任意四个非负测度的表示擅自使用这个支配关系。

反向若四个典范分量 Radon，则 $|\mu|\leqslant\sum_{j=1}^4\mu_j$，可由全变差三角不等式得到。各 $K_j\subseteq E\subseteq U_j$，故有限并 $K$ 紧、有限交 $U$ 开，且 $U\setminus K\subseteq U_j\setminus K_j$ 对每个 $j$ 成立。于是
\begin{align*}
|\mu|(U\setminus K)\leqslant\sum_{j=1}^4\mu_j(U\setminus K)
\leqslant\sum_{j=1}^4\mu_j(U_j\setminus K_j)<\varepsilon.
\end{align*}
这也说明原式必须是 $|\mu|$ 的质量，而不是复数 $\mu(U\setminus K)$ 与实数比较。对于前段由四个\emph{任意}有限正 Radon 测度产生的表示，仍有 $|\mu|\leqslant\sum_j\mu_j$，以上反向论证完全适用，所以确实得到复 Radon 测度。

更一般地，若 $0\leqslant\lambda\leqslant\rho$ 且 $\rho$ 是有限 Radon 测度，则 $\rho$ 的夹层同样控制 $\lambda$，故 $\lambda$ Radon。有限个有限正 Radon 测度之和也 Radon，证明仍取各紧集的并及各开集的交并分配误差。对 $\mu,\nu\in M(X)$ 及 $a,b\in\mathbb C$，利用
\begin{align*}
|a\mu+b\nu|\leqslant|a|\,|\mu|+|b|\,|\nu|
\end{align*}
便知其全变差 Radon，进而 $a\mu+b\nu\in M(X)$。最后，$\|\mu\|\geqslant0$；若为零，则 $|\mu(E)|\leqslant|\mu|(X)=0$ 对每个 Borel 集成立，故 $\mu=0$；反向显然。恒等式 $|a\mu|=|a|\,|\mu|$ 给出绝对齐次性，三角不等式由全变差的不等式对 $X$ 取值得到。因此它确为范数。动机是用全变差同时控制实、虚及正、负贡献，消除抵消造成的假象。
\end{remark}
% END SUPPLEMENT: id=v14-ch07-variation-radon-details
\subfile{section-03-riesz.tex}
\end{document}
